\documentclass[12pt,a4paper,oneside]{article}
\usepackage[brazil]{babel}
\usepackage[utf8]{inputenc}
\usepackage{olymp}

\setcounter{problem}{1}

\begin{document}

\begin{problem}{Deu zica}{zica}{2}
 
Pesquisadores da Universidade Chun-Li desenvolveram um poderoso veneno que prometia acabar com dengue, zika e chicungunha. O combate ao mosquito \textit{aedes aegypti} foi um sucesso... para o mosquito! Uma população inicial de mosquitos que teve contato com o veneno agora consegue se reproduzir sem necessidade de água parada e nem de outros mosquitos! Ou seja, deu zica...

Os estudos mostram que em certa região, a cada hora cada mosquito consegue gerar outros 2. Assim, partindo de apenas 1 mosquito infectado, em 1 hora serão 3 mosquitos (ele e mais dois); em 2 horas já serão 9; em 3 horas 27, e assim por diante. É possível que rapidamente a população chegue a 1.000.000.000 de mosquitos!

Faça um programa para, sabendo o tamanho da população inicial de mosquitos que teve contato com o veneno e a força do veneno, que indica quantos mosquitos cada um consegue gerar por hora, calcule em quanto tempo a população ultrapassa 1.000.000.000 de mosquitos.

%\Notes
%
%Observações, se tiver.

\InputFile

A entrada é composta por dois inteiros, $P_0$ e $G$, que representam respectivamente a população inicial de mosquitos infectados (1, no exemplo acima) e sua taxa de reprodução, ou seja, quantos mosquitos são geradas por mosquito a cada hora (2, no exemplo acima). Restrições: $1 \leq P_0, G \leq 1000$.


\OutputFile

Seu programa deve gerar apenas uma linha de saída, contendo um único inteiro que indica quantas horas serão necessárias para a população atingir no mínimo 1.000.000.000 de mosquitos.
% \Notes
% Preste atenção aos tipos das variáveis, pois $20! \approx 2^{61}$.

\Example

\begin{example}
\exmp{
1 2
}{
19
}%
\end{example}

\begin{example}
\exmp{
1 9
}{
9
}%
\end{example}

\begin{example}
\exmp{
3 9
}{
9
}%
\end{example}

\begin{example}
\exmp{
6 13
}{
8
}%
\end{example}

%Comentários sobre o exemplo, se necessário.

\end{problem}

\end{document}
